\documentclass[12pt,a4paper]{article}
\usepackage[margin=25mm, headheight=15pt, headsep=10pt]{geometry}
\usepackage{fontspec}
\usepackage[table]{xcolor} % Note: table option is required for rowcolors
\usepackage{titlesec}
\usepackage{tocloft}
\usepackage{fancyhdr}
\usepackage{booktabs}
\usepackage{tcolorbox}
\tcbuselibrary{skins, breakable}
\usepackage[colorlinks=true, linkcolor=emerald, urlcolor=emerald, bookmarks=true]{hyperref}

\definecolor{emerald}{RGB}{0,102,51}
\definecolor{darkred}{RGB}{139,0,0}
\definecolor{lightgray}{RGB}{245,245,245}

\usepackage[nil]{babel}
\babelprovide[import, main]{bengali}
\babelprovide[import]{arabic}
\babelfont{rm}[Renderer=HarfBuzz,Scale=1.0]{Noto Serif Bengali}
\babelfont{sf}[Renderer=HarfBuzz]{Noto Sans Bengali}
\babelfont{tt}[Renderer=HarfBuzz]{Noto Sans Bengali}
\babelfont[arabic]{rm}[Renderer=HarfBuzz,Scale=1.1]{Noto Naskh Arabic}

% Section styling
\titleformat{\section}{\Large\bfseries\color{emerald}}{\thesection.}{0.5em}{}
\titleformat{\subsection}{\large\bfseries\color{emerald!85!black}}{\thesubsection.}{0.5em}{}
\titleformat{\subsubsection}{\normalsize\bfseries\color{emerald!70!black}}{\thesubsubsection.}{0.5em}{}
\titlespacing*{\section}{0pt}{18pt}{8pt}

% Fancy Header/Footer setup
\pagestyle{fancy}
\fancyhf{} % clear all header and footer fields
\fancyhead[L]{\color{emerald}\small\textbf{\nouppercase{\leftmark}}}
\fancyfoot[L]{\scriptsize\color{black!50}{\lat{AI}}-সংকলিত, আলেম-যাচাইকৃত নয়}
\fancyfoot[C]{\color{emerald!70!black}\thepage}
\renewcommand{\headrulewidth}{0.4pt}
\renewcommand{\headrule}{\hbox to\headwidth{\color{emerald}\leaders\hrule height \headrulewidth\hfill}}
\renewcommand{\footrulewidth}{0pt}

% Custom boxes
% 1. Ayat/Hadith Box
\newtcolorbox{ayatbox}[1][]{
    enhanced, breakable, 
    colback=emerald!5, 
    colframe=emerald, 
    boxrule=0.5pt, 
    arc=3pt, 
    left=10pt, right=10pt, top=10pt, bottom=10pt,
    #1
}

% 2. Quote/Translation Box
\newtcolorbox{quotebox}[1][]{
    enhanced, breakable,
    frame hidden,
    colback=lightgray,
    borderline west={3pt}{0pt}{emerald},
    left=15pt, right=10pt, top=10pt, bottom=10pt,
    sharp corners,
    #1
}

% 3. AI-generated notice box
\newtcolorbox{warnbox}[1][]{
    enhanced, breakable,
    colback=red!4,
    colframe=darkred,
    boxrule=0.8pt,
    arc=3pt,
    left=10pt, right=10pt, top=10pt, bottom=10pt,
    #1
}

% Commands
\newcommand{\arab}[1]{\foreignlanguage{arabic}{#1}}
\newcommand{\kib}[1]{\textcolor{emerald}{\textbf{#1}}}

% Latin (English) fallback: Noto Serif Bengali has NO Latin glyphs
\newfontfamily\latfont[Renderer=HarfBuzz]{Noto Serif}
\newcommand{\lat}[1]{{\latfont #1}}

\setlength{\parskip}{5pt}
\setlength{\parindent}{0pt}

% Fix for missing bullet in Noto Serif Bengali
\renewcommand{\labelitemi}{$\bullet$}



\begin{document}
\begin{center}{\Large\bfseries\color{emerald} ঘটনা ১: আবু যার ও বিলাল (রাঃ) — বর্ণবাদী অহংকারের অনুশাসন}\end{center}
\vspace{1em}
\section*{ঘটনা ১: আবু যার ও বিলাল (রাঃ) — বর্ণবাদী অহংকারের অনুশাসন}

\begin{warnbox}
 \textbf{গুরুত্বপূর্ণ নোটিশ (\lat{Important} \lat{Notice}):} এই কন্টেন্টটি \lat{AI} (কৃত্রিম বুদ্ধিমত্তা) দিয়ে প্রস্তুত ও সংকলিত। \lat{AI} কঠোর সূত্র-মান নীতি মেনে শুধু নির্ভরযোগ্য উৎস থেকে তথ্য সংগ্রহ করেছে — কেবল সহীহ/হাসান হাদিস ও সহীহ সনদের আসার রাখা হয়েছে, দুর্বল সনদ বাদ দেওয়া হয়েছে। তবে এটি কোনো আলেম বা মুহাদ্দিস কর্তৃক যাচাইকৃত নয়।
\end{warnbox}


\medskip\hrule\medskip

\subsection*{১. সূত্র কার্ড}

\begin{table}[h]\centering\small
\begin{tabular}{ll}
বিষয় & বিবরণ \\
\hline
\textbf{বর্ণনাকারী} & মা'রূর ইবনে সুয়ায়দ (রহ.), আবু যার (রাঃ) থেকে \\
\textbf{গ্রন্থ} & সহীহ বুখারি, হাদিস ৩০ ও ৬০৫০ \\
\textbf{অধ্যায়} & কিতাবুল ইমান, বাব: মুসলিমরাই ভাই; কিতাবুল মনাকিব, বাব: বিলালের মর্যাদা \\
\textbf{সনদের মান} & \textbf{সহীহ} (ইমাম বুখারির শর্তানুযায়ী) \\
\textbf{মূল বিষয়} & বর্ণ, গোত্র বা সামাজিক অবস্থান নিয়ে অহংকার করা নিষিদ্ধ — এটি জাহিলিয়াতের চিহ্ন \\
\hline
\end{tabular}
\end{table}

\medskip\hrule\medskip

\subsection*{২. পটভূমি (কে, কখন, কোথায়, কেন)}

মদীনায় নববী রাষ্ট্র প্রতিষ্ঠার পর ইসলাম মানুষের মধ্যে গোত্রীয় শ্রেষ্ঠত্বের \lat{concept} (ধারণা) নির্মূল করছিল। কিন্তু দীর্ঘ জাহিলিয়াতের \lat{habit} (অভ্যাস) মানুষের মনে গেঁথে ছিল।

একবার \textbf{আবু যার গিফারী (রাঃ)} — যিনি ইসলামের প্রথম যুগে সত্যের সন্ধানে মক্কায় এসে ইসলাম গ্রহণকারী প্রসিদ্ধ সাহাবি — এবং \textbf{বিলাল ইবনে রাবাহ (রাঃ)} — যিনি ইসলামের প্রথম মুয়াযযিন, আবিসিনিয়ার কৃষ্ণাঙ্গ \lat{slave} (দাস) থেকে মুক্ত হয়ে নবীজির ঘনিষ্ঠ \lat{companion} (সাহাবি) হয়েছিলেন — তাদের মধ্যে কোনো বিষয়ে \lat{argument} (কথা কাটাকাটি) হয়।

রাগের মাথায় আবু যার (রাঃ) বিলাল (রাঃ)-কে তার \textbf{মায়ের \lat{complexion} (বর্ণ)} টেনে \lat{humiliate} (অপমান) করলেন — \textbf{\lat{“}হে কালো মায়ের সন্তান!\lat{”}}

এই এক লাইনের কথাই ঘটনার \lat{centre} (কেন্দ্র)। কেন এই ঘটনা \lat{important} (গুরুত্বপূর্ণ)? কারণ দুইজন সেরা সাহাবির মাঝেও বর্ণবাদী \lat{arrogance} (অহংকার) এর এই ছাপ ছিল — এবং নবীজি (সা.) তাৎক্ষণিকভাবে তা নির্মূল করে দিলেন।

\medskip\hrule\medskip

\subsection*{৩. পূর্ণ ঘটনার বর্ণনা}

মা'রূর ইবনে সুয়ায়দ (রহ.) বলেন:

\begin{quote}
\textbf{\lat{“}আমি 'রাবাযা' নামক স্থানে আবু যার (রাঃ)-এর সঙ্গে দেখা করলাম। তাঁর পরনে ছিল এক জোড়া কাপড় (লুঙ্গি ও চাদর), আর তাঁর ভৃত্যের পরনেও ছিল ঠিক একই ধরনের এক জোড়া কাপড়।} \\
\textbf{আমি তাঁকে এর কারণ জিজ্ঞেস করলাম। তিনি বললেন: "একবার আমি জনৈক ব্যক্তিকে গালি দিয়েছিলাম এবং তাকে তার মা সম্পর্কে লজ্জা দিয়েছিলাম (মায়ের রঙের কথা মনে করিয়ে দিয়েছিলাম)। তখন আল্লাহর রাসূল (সাল্লাল্লাহু আলাইহি ওয়া সাল্লাম) আমাকে বললেন:} \\
\textbf{'হে আবু যার! তুমি তাকে তার মা সম্পর্কে লজ্জা দিয়েছ? তুমি তো এমন ব্যক্তি, তোমার মধ্যে এখনো জাহিলিয়াতের (অজ্ঞতার যুগের) স্বভাব রয়ে গেছে।'} \\
\textbf{'জেনে রাখো, তোমাদের দাস-দাসী তোমাদের ভাই। আল্লাহ তাদের তোমাদের অধীনস্থ করেছেন। তাই যার ভাই তার অধীনে থাকবে, সে যেন নিজে যা খায় তাকে তা-ই খাওয়ায়, আর নিজে যা পরিধান করে তাকেও তা-ই পরায়। তাদের উপর এমন কাজ চাপিয়ে দিও না যা তাদের সামর্থ্যের বাইরে। যদি এমন কঠিন কাজ দিতে হয়, তাহলে তোমরা নিজেরাও তাদের সেই কাজে সাহায্য করো।'"}
\end{quote}

ঘটনার পরপরই আবু যার (রাঃ) অনুতপ্ত হন এবং বিলাল (রাঃ)-এর কাছে ক্ষমা চান। নবীজির এই শিক্ষা তার মনে এত গভীরভাবে গেঁথে গিয়েছিল যে — \textbf{রাবাযায় (তার জীবনের শেষ সময়ে) তাঁর ভৃত্যের পরনেও ছিল তাঁর নিজের মতো একই জোড়া কাপড়} — যেন অহংকারের সামান্যতম ভাবও না থাকে।

\medskip\hrule\medskip

\subsection*{৪. শব্দের ব্যাখ্যা}

\begin{table}[h]\centering\small
\begin{tabular}{ll}
শব্দ & অর্থ \\
\hline
\textbf{জাহিলিয়াত} & ইসলাম পূর্ববর্তী অজ্ঞতার যুগ — যেখানে বর্ণ-গোত্রের আধিপত্য ছিল \\
\textbf{বর্ণবাদী অহংকার} & কাউকে তার চামড়ার রং, বংশ বা শ্রেণির কারণে হেয় জ্ঞান করা \\
\textbf{ইখওয়ান / ভাই} & ইসলামে সব মুমিন পরস্পর ভাই — কোনো শারীরিক বা সামাজিক মাপকাঠিতে ভেদ নেই \\
\textbf{মা'রূর} & এই হাদিসের বর্ণনাকারী তাবেয়ি, যিনি রাবাযায় আবু যারের সাথে সাক্ষাৎ করেন \\
\hline
\end{tabular}
\end{table}

\medskip\hrule\medskip

\subsection*{৫. সংশ্লিষ্ট হাদিস ও আয়াত}

\textbf{হাদিস (বুখারি ৩০/৬০৫০):} উক্ত হাদিসের মূল বাণী — \lat{“}তোমাদের দাস-দাসী তোমাদের ভাই; তাদেরকে নিজের মতো খাওয়াও-পরাও।\lat{”}

\textbf{আয়াত (সূরা আল-হুজুরাত ৪৯:১৩):}
\begin{quote}
\textbf{\lat{“}হে মানুষ! আমি তোমাদেরকে এক পুরুষ ও এক নারী থেকে সৃষ্টি করেছি এবং জাতি-গোত্রে বিভক্ত করেছি, যাতে তোমরা পরস্পর পরিচিত হতে পারো। নিশ্চয়ই আল্লাহর কাছে তোমাদের মধ্যে সবচেয়ে মর্যাদাবান সে, যে সবচেয়ে তাকওয়াবান।\lat{”}}
\end{quote}

\textbf{হাদিস (মুসলিম ২৫৬৪):} নবীজি (সা.) বলেছেন: \textbf{\lat{“}আল্লাহ তোমাদের বর্ণ ও সম্পদ নিয়ে দেখেন না; বরং তোমাদের অন্তর ও আমল দেখেন।\lat{”}}

\textbf{হাদিস (বুখারি ৫৯৭০):} \textbf{\lat{“}যে ব্যক্তি অহংকারবশে পোশাক টেনে চলে, আল্লাহ কিয়ামতের দিন তার দিকে তাকাবেন না।\lat{”}}

\medskip\hrule\medskip

\subsection*{৬. রেফারেন্স}

\begin{itemize}
\item সহীহ বুখারি, হাদিস ৩০ (কিতাবুল ইমান, বাব: মুসলিমরাই ভাই)
\item সহীহ বুখারি, হাদিস ৬০৫০ (কিতাবুল মনাকিব, বাব: বিলালের মর্যাদা)
\item সহীহ মুসলিম, হাদিস ৯১ (কিবরের সংজ্ঞা)
\item সহীহ মুসলিম, হাদিস ২৫৬৪ (আল্লাহ বর্ণ-সম্পদ নিয়ে দেখেন না)
\item সহীহ বুখারি, হাদিস ৫৯৭০ ও ৬১১৪
\item সূরা আল-হুজুরাত, আয়াত ১৩
\end{itemize}

\end{document}
